\section{Pairing parity and a conditional rank-eight construction}
\label{sec:pairing-morita-20261027}
The two rank-six rings remain unresolved. The following obstruction concerns
actual module categories, not just nonnegative integral representations.

\begin{lemma}[Morita pairing parity]
Let $\mathcal C$ be pseudo-unitary and $A$ a connected separable algebra.
Give its module category the positive trace and $A$ the compatible symmetric
separable Frobenius structure. If a selfdual simple $Q$ of
${}_A\mathcal C_A$ has second indicator $\epsilon$, every selfdual
constituent of its underlying object having indicator $-\epsilon$ occurs
with even multiplicity.
\end{lemma}
\begin{proof}
The forgetful functor is pivotal Frobenius monoidal. It transports the
chosen nondegenerate selfduality pairing, including its pivotal transpose
sign. On an isotypic block of opposite sign, the multiplicity space must
carry a nondegenerate alternating form; its dimension is even.
The positive normalization is essential:
$\dim_{{}_A\mathcal C_A}Q=\dim_{\mathcal C}U(Q)/\dim_{\mathcal C}A$.
Thus the induced spherical structure is the canonical positive one.
The categorical inputs are Schaumann, \emph{Traces on module categories
over fusion categories}, Proposition 5.8 and Theorem 6.6
(\url{https://arxiv.org/abs/1206.5716}), and Mulevi\v cius,
\emph{Condensation inversion and Witt equivalence via generalised orbifolds},
Example 4.18 (\url{https://arxiv.org/abs/2206.02611}).
This transports pairings; it does not assert preservation of numerical
Frobenius--Schur indicators by the forgetful functor.
\end{proof}

For $\RR{8}{1660}$ the source signs on $(1,g,X,gX,\sigma,T)$ are
$(1,1,1,1,-1,1)$. The mixed-sign type-III dual contains a selfdual simple
whose restriction is $\sigma+T$. The second mixed type-IV realization
contains one restricting to $X+gX+\sigma+T$. Either sign for that simple
would leave an odd opposite-sign block. Both module realizations are
therefore impossible. The other mixed type-IV realization passes this test.
The restrictions follow by internal-Hom reciprocity:
$[U(Q):Z]=(N_ZQ)_{mm}$, using the complete dual action matrices, not a
dimension match.

\begin{proposition}[The $1+X$ construction target]
In a pseudo-unitary categorification of $\RR{8}{1660}$, any connected
separable algebra with underlying object $1\oplus X$ has bimodule category
with based ring $R8\_3\_941$. The same holds for $1\oplus gX$.
\end{proposition}
\begin{proof}
The distinguished free module is simple and has internal end $1+X$.
Its zero diagonal $\sigma$ coefficient makes it an Ising leaf.
Only type III in the complete module list has this internal end; types II
and IV have nonzero $T$ coefficient, while type I is regular. The first
diagonal sign is positive. The mixed second sign has just been excluded.
For equal signs the complete dual-action computation has two rank-eight
possibilities. The weighted second-indicator obstruction eliminates one.
The other tensor is $R8\_3\_941$: the source-to-census permutation
$(0,3,6,7,4,5,2,1)$ identifies all $512$ structure constants.
Relabelling $X,gX$ gives the other assertion.
\end{proof}

The inverse internal-end algebra in this Morita context is also $1+Y$,
with $Y$ selfdual of dimension $d=3+2\sqrt2+2\sqrt3+\sqrt6$.
Neither algebra has been constructed. Fixing the unit, the multiplication
on $X^2$ has one coefficient in $\operatorname{Hom}(X^2,1)$ and three in
$\operatorname{Hom}(X^2,X)$, before normalization and gauge. Its associativity
still depends on the unknown ambient associator. The six-simple connecting
module, consisting of two regular Ising blocks, is an exact construction
target, not an existence theorem. Equal global FP dimension alone would
not identify this partner: $R8\_3\_928$ has the same dimension but is
categorically Morita rigid.

For the negative spin representation, $(44,2)$ is impossible in low frame
zero: phase $12$ has a unique symmetric restriction column and hence no
negative charge-conjugation line there. In frame one the total asymmetric
column budget is eight. Phase-$36$ asymmetric columns cannot be discarded;
both overall spin alternatives remain. Fusion by the distinguished
Galois-unit object of dimension $d^2$ acts on the negative conjugation
space as a signed permutation: its integral symmetric matrix squares to
the identity, so every row has one entry $\pm1$.

\subsection{A forced algebra for $\RR{4}{155}$}
Write $X^*=Y$, $\dim X=\dim Y=d_0=4+3\sqrt2$.
For $A_k=1+kX+kY$, subtraction of the four distinct pointed free modules
from its induction Gram leaves vectors $a,b$ satisfying
\[
 \lVert a\rVert^2=\lVert b\rVert^2=1+8k-4k^2,
 \qquad a\cdot b=8k-4k^2.
\]
Thus $k\le2$. For $k=1$ the viable residual modules have underlying
objects $2X+Y,X+2Y$. The other Gram possibilities contain a module
underlying $X$; its internal-end dimension
$d_0^2/(1+2d_0)=2+4\sqrt2/3$ is not an algebraic integer.
For $k=2$ the residual modules underlie $X,Y$. Both surviving necessary
based modules have rank six.

The canonical central algebra for the pointed subcategory is $1+B$,
where $B$ is the Galois-unit object. Its forgetful algebra has augmentation
and splits as $1\times A_2$ with $[A_2]=1+2X+2Y$.
This forces an actual algebra in a hypothetical category. However the
residual $X$ module has internal end $1+h$, returning to the known pointed
Morita module category. It neither constructs a source category nor gives
a contradiction; existence of $A_1$ is not forced.

Complete matrices, tensor identities, audits and short finite verifiers are
in \filename{evidence/frontier_20261027/README.md}. They retain the separate
categorical premises above and do not promote necessary data to models.
